% =========================================================================
% Revised Section V: Nonlinear Direction-Finding Baseline
% Replaces the original Section "Nonlinear Estimator" in IEEE_TCOM_RV1.tex
% 
% Key changes vs original:
%   1. Reframed as "baseline" (not main contribution)
%   2. Formulated as profile MLE for direction with eta as nuisance
%   3. n_r-independent (eta absorbs cos(psi))
%   4. Explicit sphere constraint
%   5. Clarifies relationship to MLE and why GLS is preferred
% =========================================================================
\documentclass[journal]{IEEEtran}
\usepackage{amsmath,amsfonts,amssymb}
\usepackage{bm}

\begin{document}


% *******************************************
\section{Nonlinear Direction-Finding Baseline}
\label{sec:NoLineal}
% *******************************************

Before developing the closed-form estimators that constitute the main 
contribution of this work, we present an iterative nonlinear least-squares 
(NLS) approach that serves as a baseline for comparison.  This method 
directly fits the Lambertian power model to the $K$ direction-finding 
measurements and represents the class of numerical estimators commonly 
employed in optical positioning~\cite{Qin2020,Ma2024}.

\subsection{Direction-Finding Observation Model}

From \eqref{eq:Pr_i}--\eqref{eq:cosines}, the noise-free mean power 
at orientation~$i$ can be written as:
\begin{equation}
  \mu_i = \eta\,\bigl(\mathbf{n}_{t,i}\cdot\mathbf{n}_d\bigr)^{m},
  \qquad i=1,\dots,K,
  \label{eq:mu_eta}
\end{equation}
where
\begin{equation}
  \eta \triangleq \frac{C\cos\psi}{d^{2}}
  \label{eq:eta_def}
\end{equation}
is a scalar \emph{nuisance amplitude} that absorbs the unknown distance 
$d$ and the receiver-orientation factor 
$\cos\psi = \mathbf{n}_r\!\cdot\!(-\mathbf{n}_d)$.  
The only unknowns in~\eqref{eq:mu_eta} are the direction 
$\mathbf{n}_d\!\in\!\mathbb{S}^{2}$ and the scalar~$\eta>0$; 
crucially, the model does \emph{not} require knowledge of $\mathbf{n}_r$.

\subsection{NLS Formulation}

Given the sample means $\hat\mu_i = \bar P_{r,i}$ from~\eqref{eq:sampleMean}, 
the NLS direction estimator jointly minimises the sum of squared residuals 
over $\mathbf{n}_d$ and~$\eta$:
\begin{equation}
  \bigl(\hat{\mathbf{n}}_{d},\,\hat\eta\bigr)
  = \arg\min_{\substack{\mathbf{n}_d\in\mathbb{S}^{2},\;\eta>0 \\
                         \mathbf{n}_{t,i}\cdot\mathbf{n}_d\ge0,\;i=1,\dots,K}}
    \sum_{i=1}^{K}
    \bigl(\hat\mu_i - \eta\,(\mathbf{n}_{t,i}\cdot\mathbf{n}_d)^{m}\bigr)^{2}.
  \label{eq:nls_cost}
\end{equation}
The constraint $\mathbf{n}_{t,i}\cdot\mathbf{n}_d\ge0$ enforces 
$\cos\phi_i\ge0$ (LED illuminates the receiver), and the unit-sphere 
constraint $\|\mathbf{n}_d\|=1$ restricts the search to directions.

Under the i.i.d.\ Gaussian noise model of Section~\ref{sec:SystemModel}, 
minimising~\eqref{eq:nls_cost} is equivalent to maximising the 
\emph{profile likelihood} of~$\mathbf{n}_d$ after concentrating out 
the nuisance~$\eta$.  For each candidate direction, the optimal nuisance 
parameter is the least-squares solution:
\begin{equation}
  \hat\eta(\mathbf{n}_d)
  = \frac{\displaystyle\sum_{i=1}^{K}\hat\mu_i\,
    (\mathbf{n}_{t,i}\cdot\mathbf{n}_d)^{m}}
   {\displaystyle\sum_{i=1}^{K}
    (\mathbf{n}_{t,i}\cdot\mathbf{n}_d)^{2m}}.
  \label{eq:eta_hat}
\end{equation}
Hence, the NLS is the \emph{profile maximum-likelihood estimator (MLE)} 
for the direction subproblem with~$\eta$ as nuisance.

In practice, \eqref{eq:nls_cost} is solved via MATLAB's \texttt{fmincon} 
with the SQP algorithm.  The initial direction is set to the LED 
orientation~$\mathbf{n}_{t,i^*}$ corresponding to the highest received 
power ($i^* = \arg\max_i \hat\mu_i$), and $\eta$ is initialised to unity.

\subsection{Comparison with Linear Estimators}

Two observations motivate the closed-form GLS/WLS estimators developed 
in Section~\ref{sec:Lineal}:
\begin{enumerate}[leftmargin=*]
\item \textbf{Statistical weighting.}  
  The NLS cost~\eqref{eq:nls_cost} assigns equal weight to every 
  orientation, regardless of its SNR.  Orientations with negligible 
  received power ($\cos\phi_i\approx0$) contribute noise-dominated 
  residuals that degrade the solution.  By contrast, the GLS 
  estimator (Section~\ref{sec:Lineal}) applies the inverse-covariance 
  weighting $\mathbf{\Sigma}_\beta^{-1}$, which naturally down-weights 
  low-SNR orientations.

\item \textbf{Computational cost.}  
  The NLS requires an iterative constrained optimisation at each receiver 
  position, yielding ms-level latency per estimate.  The GLS/WLS 
  reduce direction finding to a single $3\times3$ eigenvalue problem, 
  achieving $\mu$s-level latency (see Section~\ref{sec:results}).
\end{enumerate}

\noindent
Despite being the profile MLE for the direction subproblem, the NLS 
does not outperform GLS in the operating regime of this work ($N=1000$ 
samples, SNR$\approx14\,$dB), because the first-order linearisation 
underlying GLS is tight and the optimal weighting 
$\mathbf{\Sigma}_\beta^{-1}$ compensates for the approximation.
The NLS therefore serves as a computationally expensive benchmark 
that validates the effectiveness of the proposed closed-form approach.






% *******************************************
\section{Nonlinear Direction-Finding Benchmark: Exact Maximum Likelihood}
\label{sec:NoLineal}
% *******************************************

Before developing the closed-form estimators that constitute the main contribution of this work, we rigorously formulate the exact Maximum Likelihood Estimator (MLE) for the direction-finding stage. While traditional numerical approaches often apply generic least-squares fitting \cite{Qin2020, Ma2024}, formulating the exact MLE establishes the absolute theoretical performance ceiling for the unapproximated nonlinear measurement model, serving as a benchmark for our proposed linear methods.

\subsection{Likelihood Function and Nuisance Parameter}

During the direction-finding stage, the $K$ beam-steered measurements are acquired from a fixed receiver position. Consequently, the absolute distance $d$ and the receiver's incidence angle $\cos\psi = \mathbf{n}_r \cdot (-\mathbf{n}_d)$ remain strictly constant across all $i \in \{1, \dots, K\}$. From the received power model in \eqref{eq:Pr_i} and \eqref{eq:cosines}, the noise-free mean power for the $i$-th orientation can be expressed exclusively in terms of the angular geometry:
\begin{equation}
    \mu_i(\mathbf{n}_d, \eta) = \eta\,\bigl(\mathbf{n}_{t,i}\cdot\mathbf{n}_d\bigr)^{m},
    \label{eq:mu_eta}
\end{equation}
where we define the scalar nuisance parameter:
\begin{equation}
    \eta \triangleq \frac{C\cos\psi}{d^{2}}.
    \label{eq:eta_def}
\end{equation}
This parameterization elegantly decouples the purely angular unknown $\mathbf{n}_d \in \mathbb{S}^{2}$ from the channel attenuation and receiver orientation. The observation vector $\mathbf{P}_r = [P_{r,1}, \dots, P_{r,K}]^{\mathsf{T}}$ follows a multivariate Gaussian distribution. Therefore, the log-likelihood function for the unknown parameters $\boldsymbol{\theta} = [\mathbf{n}_d^{\mathsf{T}}, \eta]^{\mathsf{T}}$ is given by:
\begin{equation}
    \ell(\mathbf{n}_d, \eta) = \kappa - \frac{N}{2\sigma^2} \sum_{i=1}^{K} \left( P_{r,i} - \eta \bigl(\mathbf{n}_{t,i}\cdot\mathbf{n}_d\bigr)^{m} \right)^2,
    \label{eq:log_likelihood}
\end{equation}
where $\kappa$ is a constant independent of $\boldsymbol{\theta}$.

\subsection{Profile MLE Formulation}

The true MLE is obtained by maximizing \eqref{eq:log_likelihood}. Given the homoscedastic nature of the AWGN, maximizing the log-likelihood is mathematically equivalent to minimizing the residual sum of squares subject to the physical constraints of the system. We formulate this constrained nonlinear optimization problem as:
\begin{equation}
  \bigl(\hat{\mathbf{n}}_{d,\text{NL}},\,\hat\eta\bigr)
  = \arg\min_{\substack{\mathbf{n}_d\in\mathbb{S}^{2},\;\eta>0 \\
                         \mathbf{n}_{t,i}\cdot\mathbf{n}_d\ge0}}
    \sum_{i=1}^{K}
    \bigl(P_{r,i} - \eta\,(\mathbf{n}_{t,i}\cdot\mathbf{n}_d)^{m}\bigr)^{2}.
  \label{eq:nls_cost}
\end{equation}

The strict equality constraint $\|\mathbf{n}_d\|^2 = 1$ ensures the vector remains on the unit sphere, while the inequality constraint $\mathbf{n}_{t,i}\cdot\mathbf{n}_d\ge0$ enforces the field-of-view physical boundary (i.e., the receiver cannot be illuminated from behind the LED plane). 

This formulation allows us to concentrate out the nuisance parameter. By setting the partial derivative $\partial \ell / \partial \eta = 0$, the optimal scale factor for any candidate direction $\mathbf{n}_d$ admits a closed-form least-squares solution:
\begin{equation}
  \hat\eta(\mathbf{n}_d)
  = \frac{\displaystyle\sum_{i=1}^{K} P_{r,i}\,
    (\mathbf{n}_{t,i}\cdot\mathbf{n}_d)^{m}}
   {\displaystyle\sum_{i=1}^{K}
    (\mathbf{n}_{t,i}\cdot\mathbf{n}_d)^{2m}}.
  \label{eq:eta_hat}
\end{equation}
Substituting \eqref{eq:eta_hat} back into \eqref{eq:nls_cost} yields the profile MLE, which depends solely on the 3D direction vector.

\subsection{Algorithmic Implementation and Trade-offs}

Because the objective function in \eqref{eq:nls_cost} is highly non-convex, iterative numerical solvers such as Sequential Quadratic Programming (SQP) are susceptible to local minima if initialized blindly. To guarantee convergence, we initialize the search space using a data-driven heuristic: the initial direction $\mathbf{n}_{d}^{(0)}$ is set to the orientation vector of the LED that recorded the maximum received power, i.e., $\mathbf{n}_{d}^{(0)} = \mathbf{n}_{t, i^*}$ where $i^* = \arg\max_i P_{r,i}$.

While this NL-MLE establishes the optimal accuracy bound, its reliance on iterative constrained optimization yields millisecond-level inference latency, scaling poorly for real-time tracking of highly dynamic targets. Furthermore, the unweighted nature of \eqref{eq:nls_cost} applies equal influence to all equations, whereas orientations with glancing incidence angles suffer from lower SNR. These structural and computational limitations motivate the derivation of the statistically efficient, closed-form estimators presented in the following section.







% *******************************************
\section{Nonlinear Direction-Finding Benchmark: Exact Maximum Likelihood}
\label{sec:NoLineal}
% *******************************************

Before developing the closed-form estimators that constitute the main contribution of this work, we rigorously formulate the exact Maximum Likelihood Estimator (MLE) for the direction-finding stage. While traditional numerical approaches often apply generic least-squares fitting \cite{Qin2020, Ma2024}, formulating the exact MLE establishes the absolute theoretical performance ceiling for the unapproximated nonlinear measurement model, serving as a benchmark for our proposed linear methods.

\subsection{Likelihood Function and Nuisance Parameter}

During the direction-finding stage, the $K$ beam-steered measurements are acquired from a fixed receiver position. Consequently, the absolute distance $d$ and the receiver's incidence angle $\cos\psi = \mathbf{n}_r \cdot (-\mathbf{n}_d)$ remain strictly constant across all $i \in \{1, \dots, K\}$. From the received power model in \eqref{eq:Pr_i} and \eqref{eq:cosines}, the noise-free mean power for the $i$-th orientation can be expressed exclusively in terms of the angular geometry:
\begin{equation}
    \mu_i(\mathbf{n}_d, \eta) = \eta\,\bigl(\mathbf{n}_{t,i}\cdot\mathbf{n}_d\bigr)^{m},
    \label{eq:mu_eta}
\end{equation}
where we define the scalar nuisance parameter:
\begin{equation}
    \eta \triangleq \frac{C\cos\psi}{d^{2}}.
    \label{eq:eta_def}
\end{equation}
This parameterization elegantly decouples the purely angular unknown $\mathbf{n}_d \in \mathbb{S}^{2}$ from the channel attenuation and receiver orientation. The observation vector $\mathbf{P}_r = [P_{r,1}, \dots, P_{r,K}]^{\mathsf{T}}$ follows a multivariate Gaussian distribution. Therefore, the log-likelihood function for the unknown parameters $\boldsymbol{\theta} = [\mathbf{n}_d^{\mathsf{T}}, \eta]^{\mathsf{T}}$ is given by:
\begin{equation}
    \ell(\mathbf{n}_d, \eta) = \kappa - \frac{N}{2\sigma^2} \sum_{i=1}^{K} \left( P_{r,i} - \eta \bigl(\mathbf{n}_{t,i}\cdot\mathbf{n}_d\bigr)^{m} \right)^2,
    \label{eq:log_likelihood}
\end{equation}
where $\kappa$ is a constant independent of $\boldsymbol{\theta}$.

\subsection{Profile MLE Formulation}

The true MLE is obtained by maximizing \eqref{eq:log_likelihood}. Given the homoscedastic nature of the AWGN, maximizing the log-likelihood is mathematically equivalent to minimizing the residual sum of squares subject to the physical constraints of the system. Let $F(\mathbf{n}_d, \eta)$ denote the objective cost function to be minimized:
\begin{equation}
  F(\mathbf{n}_d, \eta) = \sum_{i=1}^{K} \left( P_{r,i} - \eta\,(\mathbf{n}_{t,i}\cdot\mathbf{n}_d)^{m} \right)^{2}.
  \label{eq:nls_cost}
\end{equation}

To simplify the optimization space, we can concentrate out the nuisance parameter $\eta$ by expressing it as a function of $\mathbf{n}_d$. For any given candidate direction $\mathbf{n}_d$, the optimal scale factor $\hat{\eta}$ must satisfy the first-order necessary condition $\partial F / \partial \eta = 0$. Applying the chain rule to \eqref{eq:nls_cost} yields:
\begin{equation}
    \frac{\partial F}{\partial \eta} = -2 \sum_{i=1}^{K} \left( P_{r,i} - \eta\,(\mathbf{n}_{t,i}\cdot\mathbf{n}_d)^{m} \right) (\mathbf{n}_{t,i}\cdot\mathbf{n}_d)^{m} = 0.
    \label{eq:partial_derivative}
\end{equation}

Distributing the terms inside the summation, we obtain:
\begin{equation}
    \sum_{i=1}^{K} P_{r,i}\,(\mathbf{n}_{t,i}\cdot\mathbf{n}_d)^{m} = \eta \sum_{i=1}^{K} (\mathbf{n}_{t,i}\cdot\mathbf{n}_d)^{2m}.
    \label{eq:derivative_expansion}
\end{equation}

Solving \eqref{eq:derivative_expansion} for $\eta$ provides the closed-form least-squares solution for the nuisance parameter conditioned on the direction:
\begin{equation}
  \hat\eta(\mathbf{n}_d)
  = \frac{\displaystyle\sum_{i=1}^{K} P_{r,i}\,
    (\mathbf{n}_{t,i}\cdot\mathbf{n}_d)^{m}}
   {\displaystyle\sum_{i=1}^{K}
    (\mathbf{n}_{t,i}\cdot\mathbf{n}_d)^{2m}}.
  \label{eq:eta_hat}
\end{equation}

By substituting \eqref{eq:eta_hat} back into \eqref{eq:nls_cost}, the dimensionality of the optimization problem is reduced strictly to the spatial domain. The resulting profile MLE is formally defined as:
\begin{equation}
  \hat{\mathbf{n}}_{d,\text{NL}}
  = \arg\min_{\substack{\mathbf{n}_d\in\mathbb{S}^{2} \\
                         \mathbf{n}_{t,i}\cdot\mathbf{n}_d\ge0}}
    F\bigl(\mathbf{n}_d, \hat\eta(\mathbf{n}_d)\bigr),
  \label{eq:profile_mle}
\end{equation}
where the strict equality constraint $\|\mathbf{n}_d\|^2 = 1$ ensures the vector remains on the unit sphere, and the inequality constraint $\mathbf{n}_{t,i}\cdot\mathbf{n}_d\ge0$ enforces the field-of-view physical boundary.

\subsection{Algorithmic Implementation and Transition to Linear Estimators}

Because the concentrated objective function in \eqref{eq:profile_mle} is highly non-convex, iterative numerical solvers such as Sequential Quadratic Programming (SQP) are susceptible to local minima if initialized blindly. To guarantee convergence, we initialize the search space using a data-driven heuristic: the initial direction $\mathbf{n}_{d}^{(0)}$ is set to the orientation vector of the LED that recorded the maximum received power, i.e., $\mathbf{n}_{d}^{(0)} = \mathbf{n}_{t, i^*}$ where $i^* = \arg\max_i P_{r,i}$.

While this NL-MLE establishes the optimal theoretical accuracy bound by exactly modeling the nonlinear measurement process, its reliance on iterative constrained optimization yields millisecond-level inference latency, scaling poorly for real-time tracking in low-power receivers. Furthermore, the unweighted nature of the cost function applies equal geometric influence to all equations, disregarding the heteroscedastic impact of orientations with glancing incidence angles and lower signal-to-noise ratios. These structural and computational limitations motivate the derivation of the statistically efficient, closed-form generalized least squares (GLS) estimator presented in the following section.



\end{document}
